The system speaks on a user's behalf, so added content, polarity changes, and inferred intentions can misrepresent that person.
Candidate checks reduce some observed errors but are not a substitute for user control, caregiver review, or independent language assessment.
No study with AAC users is reported, and synthetic dialogues are not evidence of clinical effectiveness.
Voice conditioning should use recordings with explicit speaker permission; placeholder tones are labelled to avoid presenting them as synthesized speech.
Personal vocabulary and voice recordings can be sensitive and require appropriate consent and access controls in any deployment.
